% 本文件由 LaTeX 排版 Agent 根据有效 translation.json 自动生成。
% author=Miscellaneous; work_id=0857; title=执事哈比卜殉道记

\chapter{执事哈比卜殉道记}

% structure: p0001 source=h1 target=omitted reason=page-title-already-rendered-by-parent

在马其顿人亚历山大王国六百二十年，阿布月，利西尼乌斯和君士坦丁执政之年，即他出世之年，尤利乌斯和巴拉克任行政长官时期，厄德萨主教科纳的日子，利西尼乌斯在戴克里先皇帝进行的第一次迫害之后，对教会和所有基督徒人民发起了迫害。利西尼乌斯皇帝命令举行祭献和奠祭，并在各处修复祭台，以便在宙斯面前焚烧香料和乳香。\par

当许多人受到迫害时，他们自愿呼喊：\Quote{我们是基督徒}；他们不惧怕迫害，因为受迫害的人比那迫害\Emph{他们}的人更多。\par

哈比卜来自特尔泽哈村庄，已被任命为执事，他秘密进入各村庄的教会，履行职务并宣读圣经，用言语鼓励和坚固许多人，劝诫他们坚定持守信仰的真理，不要惧怕迫害者，并给他们指示。\par

当许多人因他的话得到坚固，并亲切地接受他的教导，小心不背弃所立的盟约时，城中的沙里尔——那些被特别委任处理此事的人——听到了这事，就去告诉在厄德萨镇的总督利萨尼亚斯，对他说：\Quote{哈比卜，特尔泽哈村庄的一位执事，到处秘密活动，抵制皇帝的命令，毫不惧怕。}\par

总督听到这些话，对哈比卜充满愤怒；他写了一份报告，派人告知利西尼乌斯皇帝哈比卜所做的一切；\Emph{他}也想了解对哈比卜以及\Emph{其余}不祭献的人会下达什么命令。\Emph{因为}虽然已经命令每个人都应祭献，但尚未命令如何处置那些不祭献的人：因为他们听说高卢和西班牙的统帅君士坦丁成了基督徒，不祭献。利西尼乌斯皇帝\Emph{如此}命令总督利萨尼亚斯：\Quote{无论何人，胆敢违抗我们的命令，我皇已下令将他用火焚烧；其他不同意祭献的人一律用刀剑处死。}\par

当这道命令传到厄德萨镇时，被举报的哈比卜已经越过\Emph{河}，到了宙格玛人的地区，也要在那里秘密服务。总督派人到他的村庄和周围各地寻找他，却找不到，他就下令逮捕他所有的家人以及村庄的居民；他们逮捕了他们，给他们戴上镣铐，包括他的母亲和其他家人，以及一些村民；并将他们带到城里，关进监狱。\par

哈比卜听说所发生的事，心中思忖，焦虑地思索：\Quote{对我而言，这样做是合宜的，}他说，\Quote{我应该去出现在审判官面前，而不是我躲藏起来，别人却因我被带到他面前，\Emph{以殉道}获得冠冕，而我自己却陷入极大的羞愧。因为逃避宣认基督信仰的人，基督信仰之名对他有何益处？看！即使他逃避此事，无论他到哪里，自然死亡都摆在他面前，他无法逃脱，因为这是对亚当所有子孙所定的律。}\par

哈比卜起身，秘密前往厄德萨，已准备好他的背部承受鞭打，两肋承受铁梳，身体承受火焚。他立刻去见老兵特奥泰克纳，他是总督随从队的首领；对他说：\Quote{我是特尔泽哈的哈比卜，就是你们寻找的人。}特奥泰克纳对他说：\Quote{如果没有人看见你来找我，请听我对你说的话，离开并回到你之前所在的地方，在这\Emph{迫害的}时期待在那里；至于你来见我，与我交谈，我这样劝你的事，不要让任何人知道或察觉。关于你的家人和村民，完全不必忧虑：不会有人伤害他们；他们只会被囚禁几天，\Emph{然后}总督就会释放他们：因为皇帝没有对他们下任何严重或令人恐惧的命令。但如果你相反不听我劝告这些事，你的血与我无干：因为如果你出现在审判官面前，根据皇帝对你下达的命令，你将无法逃脱火刑。}\par

哈比卜对特奥泰克纳说：\Quote{我所忧虑的不是我的家人和村民，而是我自己的救恩，免得它失落。我也非常苦恼，那天总督查问我时，我恰巧不在村中，因我缘故，看哪！许多人被戴上了镣铐，我被视为逃犯。因此，如果你不肯同意我的请求，把我带进总督面前，我就独自前去见他。}\par

特奥泰克纳听他说这话，就紧紧抓住他，把他交给他的助手；他们一起带他去总督的审判厅。特奥泰克纳进去禀告总督，对他说：\Quote{大人寻找的特尔泽哈的哈比卜已经来了。}总督说：\Quote{是谁带来的他？他们在哪里找到他的？他在那里做了什么？}特奥泰克纳对他说：\Quote{他本人自己来的，出于自愿，没有受到任何人的强迫，因为没有人知道他。}\par

总督听见这话，就对他极其恼怒，这样说：\Quote{这个家伙如此行事，对我极其藐视，轻视我，视我为无能的审判官；因为他这样做了，不宜对他显任何慈悲；我也不应按照皇帝关于他的命令，急于对他宣判死刑；但我应当对他忍耐，好使他受的苦刑和惩罚更为加倍，并藉着他使许多\Emph{其他}人不敢再胆敢逃跑。}\par

许多人都聚集在审判厅门口站在他身旁，有些是侍从，有些是城中百姓，其中有人对他说：\Quote{你做得不好，竟在没有审判官强迫的情况下前来，向寻找你的人显现自己。}另有人却又对他说：\Quote{你做得很好，自愿前来显现自己，而不是等审判官强迫你带来；因为现在你的信仰宣认是出于自愿，而非人的强迫。}\par

城中的\Person{沙里尔}们听见那些站在审判厅门口对他说话的人所说的话——特别是他秘密前往\Person{特奥特克纳}那里，且不愿告发他一事——也被城中的\Person{沙里尔}们听说了；他们所听见的一切都禀告了审判官。\par

审判官对那曾对\Person{哈比卜}说\Quote{你为何未经审判官强迫就前来向审判官显现自己}的人发怒。他对\Person{特奥特克纳}说：\Quote{一个人被立为同僚之首，却这样欺诈地对待上司，藐视皇帝针对叛徒\Person{哈比卜}所下、要将他用火焚烧的命令，是不合宜的。}\par

\Person{特奥特克纳}说：\Quote{我没有欺诈地对待同僚，也无意藐视皇帝所下的命令；因为在阁下面前我算什么，竟敢如此行事？但我曾严厉地审问他，为阁下所向我追问的那件事，为要知道并看清他是自愿前来，还是阁下强迫他由他人带来；当我听他说是自愿前来时，我小心地带他到大人审判厅的尊贵门口。}\par

总督急速下令，他们将\Person{哈比卜}带到他面前。差役说：\Quote{看，他站在阁下面前。}\par

他开始这样审问他，对他说：\Quote{你叫什么名字？从哪里来？你是做什么的？}\par

他对他说：\Quote{我叫\Person{哈比卜}，从\Place{特尔泽哈}村来，被立为执事。}\par

总督说：\Quote{你为什么违反皇帝的命令，执行你执事的职务——皇帝禁止你这样做——并拒绝向皇帝所敬拜的\Person{宙斯}献祭？}\par

\Person{哈比卜}说：\Quote{我们是基督徒；我们不敬拜人的作品——那些人虚无，他们的作品也虚无——但我们敬拜那创造人的天主。}\par

总督说：\Quote{不要坚持你来到我面前时的那狂妄心志，也不要侮辱\Person{宙斯}，那皇帝的伟大夸耀。}\par

\Person{哈比卜}说：\Quote{但这\Person{宙斯}是偶像，是人的作品。你说我侮辱他，这很好。然而，若把他从木头上雕刻出来、用钉子钉牢，这本身就在高声宣告他是被造的，你怎么对我说我侮辱他呢？看，他的侮辱来自于他自己，且是针对他自己的。}\par

总督说：\Quote{正是因你拒绝敬拜他，你侮辱了他。}\par

\Person{哈比卜}说：\Quote{但如果因我不敬拜他而侮辱了他，那么那用铁斧雕刻他的木匠，以及那击打他并用钉子钉牢他的铁匠，对他施加了多大的侮辱啊！}\par

总督听见他这样说，就下令毫不怜悯地鞭打他。他被五\Emph{个人}鞭打后，总督对他说：\Quote{你现在愿意服从皇帝吗？如果你不服从\Emph{他们}，我要用铁梳狠狠撕扯你，用\Emph{各种}酷刑折磨你，最后下令把你用火焚烧。}\par

\Person{哈比卜}说：\Quote{这些你用来恐吓我的威胁，比我心中已决定忍受的，还要卑劣不足道；所以我前来，出现在你面前。}\par

总督说：\Quote{把他放入杀人的铁桶里，按照他应得的鞭打他。}他被鞭打后，他们对他说：\Quote{向众神献祭。}但他大声喊道：\Quote{你们的偶像可咒骂，那些与你们一同像你们一样敬拜它们的人也可咒骂。}\par

总督下令，他们把他带到监狱；但根据审判官的命令，他们不准他与家人或同村的人说话。那天正是皇帝的节庆。\par

在\Place{伊卢尔}月第二日，总督下令，他们把他从监狱带来。他对他说：\Quote{你愿意放弃你所宣认的信仰，服从皇帝所下的命令吗？因为如果你不服从，我要用铁梳的残酷撕扯使你服从。}\par

\Person{哈比卜}说：\Quote{我没有服从他们，而且我心中已决定不服从他们——即使你加给我比皇帝所命令的更严厉的惩罚，我也不服从。}\par

总督说：\Quote{我指着众神发誓，如果你不献祭，我就不留下一件严酷痛苦的惩罚不用来折磨你；我们要看看你所敬拜的基督是否拯救你。}\par

\Person{哈比卜}说：\Quote{凡敬拜基督的人都藉着基督获得拯救，因为他们不敬拜受造物连同创造受造物的创造主。}\par

总督说：\Quote{把他拉长，用鞭子鞭打，直到他身体上没有一处未被打过。}\par

哈比卜说：\Quote{至於这些你认为因其撕裂而\Emph{如此}痛苦的折磨，从中却为忍受它们的人编织了胜利的冠冕。}\par

总督说：\Quote{你们怎能称折磨为安逸，并将身体的酷刑算作胜利的冠冕呢？}\par

哈比卜说：\Quote{这些事不是你该问我的，因为你的不信配不上听其中的缘由。我不祭祀，\Emph{我已经}说过，\Emph{现在仍然}这样说。}\par

总督说：\Quote{你受这些刑罚，是因为你该当受它们。我要挖出你的眼睛，它们曾注视这尊宙斯像而不畏惧他；我要堵住你的耳朵，它们曾听闻皇帝的法律而不战栗。}\par

哈比卜说：\Quote{你所否认的天主，那另一个世界是属于祂的；在那里，你将\Emph{被迫}以鞭打承认祂，尽管你再次否认了祂。}\par

总督说：\Quote{别提你说的那个世界了，现在焦虑地想清楚：你正受的这刑罚，没有人能救你脱离；除非众神在你向他们献祭时救你。}\par

哈比卜说：\Quote{那些为基督之名而死、不敬拜那些被制造被创造的物件的人，必在天主面前得生命；但那些贪爱短暂生命胜过那生命的——他们的痛苦将是永久的。}\par

总督下令，他们就把哈比卜吊起来，用铁梳撕裂他；撕裂的时候又推打他。他悬挂了很久，直到两臂的肩胛骨嘎吱作响。\par

总督对他说：\Quote{如今你肯顺从吗？在那边向宙斯献香？}\par

哈比卜说：\Quote{在这些痛苦之前，我没有顺从你的要求；\Emph{而}如今我已经受了它们，你怎么会以为我要顺从，从而失去我因它们所获得的一切呢？}\par

总督说：\Quote{我准备用比这些更猛烈更痛苦的刑罚叫你服从，遵照皇帝的命令，直到你遵行他们的旨意。}\par

哈比卜说：\Quote{你因我不服从皇帝的命令而惩罚我，然而，皇帝提拔你、使你成为审判官，你自己也违反了他们的命令，因为你并没有把皇帝吩咐你的事加在我身上。}\par

总督说：\Quote{因我对你忍耐，\Emph{所以}你才这样说话，像个控告的人。}\par

哈比卜说：\Quote{若你没有鞭打我、捆绑我、用铁梳撕裂我、用镣铐锁我的脚，那还\Emph{可以}认为你对我有忍耐。但既然这些事同时发生，你所说的忍耐在哪里呢？}\par

总督说：\Quote{你所说的这些对你有害无益，因为一切都对你不利，它们会给你带来甚至比皇帝所吩咐的更痛苦的折磨。}\par

哈比卜说：\Quote{我若不知道它们对我有益，就不会在你面前提起关于它们的只言片语。}\par

总督说：\Quote{\Emph{我}要使你缄默，同时平息那些你未曾敬拜的众神对你的怒气，并满足皇帝对于你违背他们命令的态度。}\par

哈比卜说：\Quote{我不怕你想用来恐吓我的死亡；我若怕死，就不会挨家挨户地去服务了：正因如此，我才那样服务。}\par

总督说：\Quote{你怎么敬拜并尊崇一个人，却拒绝敬拜并尊崇那边的宙斯呢？}\par

哈比卜说：\Quote{我并非敬拜一个人，因为圣经教导我说：\Quote{凡信赖人的人，是可咒骂的。}但那位取了身体并成为人的天主，\Emph{祂}，我才敬拜、才荣耀。}\par

总督说：\Quote{照皇帝所吩咐的去做；至于你心中所想的，若你愿意放弃，\Emph{那好}；若你不愿意，\Emph{那}就不要放弃。}\par

哈比卜说：\Quote{这两件事都做到是不可能的；因为谎言与真理相反，而那在我心中牢牢扎根的，不可能被驱除。}\par

总督说：\Quote{我要用痛苦严厉的刑罚，使你抛却你所说\InnerQuote{在我心中牢牢扎根}的念头。}\par

哈比卜说：\Quote{\Emph{至于}这些你认为会从我心中根除它的刑罚，恰恰通过这些刑罚，它在我心中生长，像一棵结果子的树。}\par

总督说：\Quote{鞭打和铁梳对你那棵树有何帮助？尤其是当我下令用火毫不怜悯地烧毁它时。}\par

哈比卜说：\Quote{我看的不是那些你眼睛看到的东西，因为我默观那看不见的事物；因此我遵行天主——\Emph{万有}的创造者——的旨意，而非遵行\Emph{人手}所做的偶像的旨意，它什么也感觉不到。}\par

总督说：\Quote{因他如此否认皇帝所敬拜的神明，要在他先前的撕裂之外再用铁梳撕裂他；因为在我耐心审问他的众多问题中，他已忘了先前的撕裂。}\par

当他们撕裂他时，他大声呼喊说：\Quote{这时的苦难，与那要显现在那些爱基督的人身上的荣耀是无法相比的。}\par

总督见他即使在如此刑罚下仍不肯献祭，就对他说：\Quote{你的教义就这样教你憎恨自己的身体吗？}\par

哈比卜说：\Quote{不，我们不憎恨自己的身体：圣经明确教导我们说：\Quote{凡丧掉生命的，必得着生命。}但它还教导我们另一件事：我们不可\Quote{把圣物给狗，也不可把珍珠丢在猪前。}}\par

总督说：\Quote{我知道你这样说话的唯一目的，就是要激怒我的怒气和我内心的愤怒，使我迅速对你宣判死刑。因此，我不会急于满足你的愿望，而是会忍耐：不是为你的解脱，而是为了使加在你身上的酷刑加重，让你亲眼看见你的肉在铁梳划过你的肋旁时从你面前掉下来。}\par

哈比布说：\Quote{我自己也期待着，你应当正如你所说的那样，在我身上加重你的酷刑。}\par

总督说：\Quote{服从皇帝们吧，他们有权随心所欲地行事。}\par

哈比布说：\Quote{随心所欲地行事不是人的本性，而是天主的本性，他的权能在天上，也统御地上的一切居民；\InnerQuote{没有人能斥责他的手，向祂说：祢在做什么？}}\par

总督说：\Quote{对于你的这种傲慢，处死剑下都太轻了。但我准备下令对你\Emph{施加}一种比剑更痛苦的死亡。}\par

哈比布说：\Quote{我也期待一种比剑更慢的死亡，你可以随时对我宣判。}\par

于是总督继续对他宣判死刑。他在侍从面前大声呼喊，在他们以及城中的贵族们聆听时说道：\Quote{这个哈比布否认了众神，你们也亲耳听到了，并且还辱骂了皇帝们，他的生命应当从这荣耀的太阳之下被涂抹，他不再\Emph{任何}应当看见这发光的星辰，众神的同伴；若非前朝皇帝们命令谋杀者的尸体应当被埋葬，这个\Emph{家伙}的尸体也不配被埋葬，因为他如此傲慢。我命令：把一条带子塞进他的嘴里，如同塞进一个谋杀者的嘴里，然后用慢火慢慢焚烧，以增加他死亡的痛苦。}\par

他从总督面前出去，嘴里塞着带子；一大群城里的人跟在他后面奔跑。基督徒们欢喜，因为他没有偏离，也没有放弃他的岗位；但外教人威胁他，因为他拒绝献祭。他们领着他经过西面的拱门，对着由\Person{Abshelama}（\Person{Abgar}的儿子）建造的墓园。他的母亲穿着白衣，与他一同出去。\par

当他到达将要焚烧他的地方时，他站起来祈祷，所有与他同来的人也祈祷；他说：\Quote{噢，基督君王，既然世界属于祢，来世也属于祢，请看并见证，我本可逃避这些苦难，却没有逃避，以免落入祢公义的手中：愿这焚烧我的火，在祢面前成为我的补偿，使我得以从那永不熄灭的火中得救；并藉着祢的圣神，接纳我的灵魂到祢面前，噢，可敬之父的荣耀圣子！}他祈祷完毕后，转身祝福了他们；他们哭泣着向他致意，无论男女；他们对他说：\Quote{请在你的主面前为我们祈祷，愿他在他的子民中赐下和平，并恢复他被倾覆的教会。}\par

当哈比布站着的时候，他们挖了一个地方，把他带来安置在里面；在他旁边立了一根木桩。他们来要把他绑在木桩上，但他对他们说：\Quote{我不会从你们将要焚烧我的这个地方挪开。}他们拿来了柴捆，摆放整齐，放在他的四周。当火焰燃起，火势猛烈升起时，他们向他喊叫：\Quote{张开你的嘴。}他一张嘴，他的灵魂就升了上去。男女众人都放声哭泣。\par

他们把他拉出来，从火中拖出，用细麻布和上等的膏油与香料覆盖他。他们抢走了他焚烧时\Emph{所用}的一些木块，弟兄们和一些平信徒抬走了他。他们为他预备安葬，将他与殉道者\Person{Guria}和\Person{Shamuna}葬在同一坟墓中，就是他们被安放在\Place{Baith Allah Cucla}的山丘上，反复为他吟诵圣咏和圣歌，深情而尊敬地把他的烧焦的身体\Emph{送到坟墓}。甚至一些犹太人和外教人也与基督徒弟兄一同参与包裹和埋葬他的身体。在他被焚烧时，以及在他被埋葬时，笼罩着场内场外的人都是一片悲伤的景象；泪水从所有人的眼中流下；每个人都光荣天主，因为他为了天主的圣名而将自己的身体交付给火焰焚烧。\par

他被焚烧的日子是\Emph{安息日}前夕，\OriginalTerm{伊鲁尔}{Ilul}月第二日——正是消息传来\Person{君士坦丁大帝}从西班牙内陆出发前往意大利城罗马，要与\Person{李锡尼}（那位现今统治罗马领土东部地区的\Emph{皇帝}）作战的日子；看哪！各地都在骚动，因为无人知道他们中谁会胜利并继续保持皇权。通过这个消息，对教会的迫害暂时缓和了一些。\par

书记官们记下了他们从法官那里听到的一切；城里的\OriginalTerm{沙里尔}{Sharirs}记下了审判厅门外所说的一切其他事情，按照既有的习俗，他们向法官报告了他们所见所闻的一切，法官的判决都记录在他们的案卷中。\par

我，\Person{德奥斐洛}，弃绝了我祖先的邪恶遗产，并宣认了基督，仔细抄写了哈比卜的这部行传，如同我先前抄写\Emph{那些}古里亚和沙穆纳的殉道者行传一样。而他曾因他们死于刀下而祝福了他们，他自己也得以相似他们，死于焚身的烈火，领受了荣冠。我记下了这些殉道者加冕的年份、月份和日子，并非为了那些像我一样亲眼目睹的人，而是为使日后的人可以知道这些殉道者在何时受苦，以及他们是怎样的人；\Emph{正如他们也可以}从早先殉道者的行传中得知，那些殉道者在\Person{多米提安}的日子以及所有同样迫害教会、并用鞭打和铁梳\Emph{撕裂}、用酷刑、利剑、火焰、可怕的大海和无情的矿坑处死了许多人的其他皇帝的日子里\Emph{受苦}。所有这一切以及类似的事，他们\Emph{都忍受了}，为了那将来赏报的希望。\par

此外，这些殉道者以及我所听闻的那些人的苦难，开了我\Person{德奥斐洛}的眼，光亮了我的心灵，我就宣认基督：祂是天主之子，并且是天主。愿这些殉道者脚上的尘土——就是他们正前去领受荣冠时，我追随在他们后面所取得的——为我求得赦免，因为我曾否认了祂；并愿祂如今在我宣认祂之后，在祂的朝拜者面前也宣认我！\par

在审判官向\Person{哈比卜}提出的第二十七个问题之后，他宣判他受火焚之刑。\par

\Emph{这里}结束执事\Person{哈比卜}殉道录。\par

\SourceRecord{Translated by B.P. Pratten.；Ante-Nicene Fathers；Vol. 8.；Edited by Alexander Roberts, James Donaldson, and A. Cleveland Coxe.；Buffalo, NY: Christian Literature Publishing Co.,；1886.；<http://www.newadvent.org/fathers/0857.htm>.}
